\documentclass[10pt,a4paper]{amsart}
\usepackage[margin=19mm]{geometry}
\usepackage{amsmath,amssymb,mathtools}
\usepackage[scr=rsfs,cal=euler]{mathalfa}
\usepackage{xfrac}
\usepackage[unicode,hidelinks,bookmarks=false]{hyperref}
\newcommand{\ie}{i.e.\ }
\newcommand{\cf}{cf.\ }
\newcommand{\resp}{respectively\ }
\newcommand{\Subsection}{\subsection}
\providecommand{\nobreakdash}{\nobreak\hbox{-}}
\setlength{\emergencystretch}{2em}
\multlinegap0pt
\usepackage{amssymb}
\usepackage[normalem]{ulem}
\usepackage{xcolor}
\usepackage[shortlabels]{enumitem}
\newtheorem{theorem}{Theorem}
\newcommand{\R}{{\mathbb R}}
\newcommand{\strt}[1]{\rule{0pt}{#1}}
\title{From generalized Poincar\'e to Poincar\'e-Sobolev inequalities via self-improving methods}
\author{Alejandro Claros, Carlos P\'erez, Linfei Zheng}
\date{Source: arXiv:2606.08556v2 (CC BY 4.0)}
\begin{document}
\maketitle

\noindent\emph{Source and licence.} From generalized Poincar\'e to Poincar\'e-Sobolev inequalities via self-improving methods. Version arXiv:2606.08556v2 (CC BY 4.0), distributed by arXiv under the Creative Commons Attribution 4.0 International licence, http://creativecommons.org/licenses/by/4.0/, as recorded in the arXiv licence metadata for this exact version and shown on the abstract page https://arxiv.org/abs/2606.08556v2. Full licence notice: \texttt{licenses/arxiv-2606.08556-CC-BY-4.0.txt}.\par\smallskip

\noindent\textit{Editorial scope.} Counted unit: the article's theorem 4 (source0.tex lines 780--786). The article's complete original proof of it is delivered with it (source0.tex lines 1718--1837, one \texttt{proof} environment of 120 lines, which the article places away from the statement). No mathematics is written, repaired or re-derived: the statement and the proof are the article's own lines, in the article's own notation, and no retained line is substituted or re-flowed.  Retained uncounted as the source-local context the proof needs: lines 367--369; lines 838--841; lines 1707--1714.  Nothing was omitted from any retained span, so the package carries no omission record.  No retained line contains a citation, so the wrapper carries no bibliography and no bibliography entry.  No retained line contains a figure environment, an image-inclusion command or drawing code, so no image is delivered and the package declares no asset dependency.  Editorial formatting only, in addition: an AMS \texttt{amsart} preamble that restates the article's own declarations and macros used by the retained lines, namely the environments \texttt{theorem} on separate counters, together with the standard packages those lines need.  The retained environments print in this document as Theorem 1 in order of appearance, whereas the article prints the same items with the numbering of its own full document.  The complete native source of the article as distributed in the arXiv e-print for this exact version is bundled unchanged as \texttt{sources/202337/source0.tex}.
\begin{equation}\label{I1 P11}
	\frac{1}{|Q|}\int_Q |f(x)-f_Q|\,dx \le c_n \,\ell(Q)\,\frac{1}{|Q|}\int_Q |\nabla f(x)|\,dx, 
\end{equation} 

\begin{equation}\label{Holder Lorentz}
	\int_{E} |f(x)g(x)|\,d\mu(x)
	\le \|f\|_{\strt{2.3ex} L^{p,q}(E,\mu)}\,\|g\|_{ \strt{2.3ex} L^{p',q'}(E,\mu)},
\end{equation}

\begin{theorem}\label{Thm FPW}
    Let $n\ge 2$ and let $\mu$ be a non-negative Borel measure. Then, there exists a dimensional constant $c_n>0$ such that for all Lipschitz functions $f$ we have
        \begin{equation}\label{Thm FPW Eq}
            \left\| f-f_Q\right\|_{\strt{2.3ex} L^{n',\infty}(Q,\mu)}
            \le c_n \int_Q |\nabla f(x)| M^c(\chi_Q\mu)(x)^\frac{1}{n'} dx,
        \end{equation}
    for each cube $Q$.
\end{theorem}

\begin{theorem}\label{Thm PS 4}
    Let $w$ be a weight in $\R^n$ with $n\ge 2$ and let $1<p<n$. Suppose $f$ is a Lipschitz function. Then, there exists a constant $c_{n,p}>0$ such that
        \begin{equation}\label{Thm PS 4 Eq}
            \left\| f-f_Q\right\|_{ \strt{2.3ex} L^{p^*,\infty}(Q,w)} \le c_{n,p} \left\| \left| \nabla f \right| \frac{M^c(w\chi_Q)^\frac{1}{n'}}{w} \right\|_{\strt{2.3ex} L^{p,\infty}(Q,w)}
        \end{equation}
    for each cube $Q$, where $\frac{1}{p^*}=\frac{1}{p}-\frac{1}{n}$ and $M^c$ denotes the centered Hardy-Littlewood maximal function.
\end{theorem}

\begin{proof}[Proof of Theorem \ref{Thm PS 4}]
By the definition of the Lorentz quasi-norm and the pointwise equivalence of the Hardy-Littlewood maximal operator $M$ and the centered Hardy-Littlewood maximal operator $M^c$, it suffices to prove that there exists $c_{n,p}>0$ such that 
%
\begin{equation}\label{Thm PS 4 normalized Eq}
    \left\| f-f_Q\right\|_{\strt{2.3ex} L^{p^*,\infty}\left(Q,\frac{w(x)dx}{w(Q)}\right)} \le c_{n,p} w(Q)^\frac{1}{n} \left\| \left| \nabla f \right| \frac{M(w\chi_Q)^\frac{1}{n'}}{w} \right\|_{\strt{2.3ex} L^{p,\infty}\left(Q,\frac{w(x)dx}{w(Q)}\right)}.
\end{equation}
For simplicity, set
\begin{equation*}
    A:=w(Q)^\frac{1}{n} \left\| \left| \nabla f \right| \frac{M(w\chi_Q)^\frac{1}{n'}}{w} \right\|_{\strt{2.3ex}L^{p,\infty}\left(Q,\frac{w(x)dx}{w(Q)}\right)}.
\end{equation*}
If $A=0$, the desired estimate is trivial. Hence, we may assume that $A>0$. Let $t_0= C_{n,p}A$, where $C_{n,p}>0$ is a sufficiently large constant to be fixed below. Since
\begin{align*}
    \left\| f-f_Q\right\|_{\strt{2.3ex}L^{p^*,\infty}\left(Q,\frac{w(x)dx}{w(Q)}\right)} = \max \bigg\{& \sup_{0<t\le t_0} t \left(\frac{w(\{x\in Q:|f(x)-f_Q|>t\})}{w(Q)}\right)^\frac{1}{p^*}, \\
    & \sup_{t_0\le t} t \left(\frac{w(\{x\in Q:|f(x)-f_Q|>t\})}{w(Q)}\right)^\frac{1}{p^*} \bigg\},
\end{align*}
we consider two cases.

\textbf{Case 1:}
\begin{equation*}
    \sup_{0<t\le t_0} t \left(\frac{w(\{x\in Q:|f(x)-f_Q|>t\})}{w(Q)}\right)^\frac{1}{p^*} \ge \sup_{t_0\le t} t \left(\frac{w(\{x\in Q:|f(x)-f_Q|>t\})}{w(Q)}\right)^\frac{1}{p^*}.
\end{equation*}
In this case, we have
\begin{equation*}
    \left\| f-f_Q\right\|_{\strt{2.3ex} L^{p^*,\infty}\left(Q,\frac{w(x)dx}{w(Q)}\right)} \le t_0=C_{n,p}A,
\end{equation*}
and the desired estimate follows.

\textbf{Case 2:}
\begin{equation*}
    \sup_{0<t\le t_0} t \left(\frac{w(\{x\in Q:|f(x)-f_Q|>t\})}{w(Q)}\right)^\frac{1}{p^*} \le \sup_{t_0\le t} t \left(\frac{w(\{x\in Q:|f(x)-f_Q|>t\})}{w(Q)}\right)^\frac{1}{p^*}.
\end{equation*}
If the last supremum is zero, then the desired estimate is immediate. Otherwise, by the definition of the supremum, there exists $t_1\ge t_0$ such that
\begin{align*}
    \frac{t_1}{2} \left( \frac{w\left(\left\{x\in Q:|f(x)-f_Q|>\frac{t_1}{2}\right\}\right)}{w(Q)} \right)^\frac{1}{p^*} 
    &\le \left\| f-f_Q\right\|_{\strt{2.3ex} L^{p^*,\infty}\left(Q,\frac{w(x)dx}{w(Q)}\right)} \\
    &= \sup_{t_0\le t} t \left(\frac{w(\{x\in Q:|f(x)-f_Q|>t\})}{w(Q)}\right)^\frac{1}{p^*} \\
    &\le 2t_1 \left( \frac{w\left(\left\{x\in Q:|f(x)-f_Q|>t_1\right\}\right)}{w(Q)} \right)^\frac{1}{p^*}.
\end{align*}
This implies that 
\begin{equation}\label{level comparison fQ}
    w\left(\left\{x\in Q:|f(x)-f_Q|>\frac{t_1}{2}\right\}\right) \le c_{n,p} w\left(\left\{x\in Q:|f(x)-f_Q|>t_1\right\}\right).
\end{equation}
Consider the set defined by 
\begin{equation*}
    E_{t_1}:=\{x\in Q:|f(x)-f_Q|>t_1\}.
\end{equation*}
We proceed by considering the truncation operator
$\tau_\lambda$ defined as follows:
\begin{equation*}
    \tau_\lambda(g)(x):=\begin{cases}0, & \text{if } g(x)\le \lambda,\\ g(x)-\lambda, & \text{if } \lambda<g(x)\le 2\lambda,\\ \lambda, & \text{if } g(x)>2\lambda.\end{cases}
\end{equation*}
Note that, for each $x\in E_{t_1}$, we have
\begin{equation*}
    \frac{t_1}{2}=\tau_{\frac{t_1}{2}}(|f-f_Q|)(x) \le \left| \tau_{\frac{t_1}{2}}(|f-f_Q|)(x)-\left(\tau_{\frac{t_1}{2}}(|f-f_Q|)\right)_Q \right|+\left(\tau_{\frac{t_1}{2}}(|f-f_Q|)\right)_Q.
\end{equation*}
Using the fact that $\tau_\lambda(g)(x)\le g(x)$ and \eqref{I1 P11}, we obtain
\begin{align*}
    \left(\tau_{\frac{t_1}{2}}(|f-f_Q|)\right)_Q
    &\le \frac{1}{|Q|}\int_Q |f(x)-f_Q|dx \\
    &\le c_n \ell(Q) \frac{1}{|Q|} \int_Q |\nabla f(x)|dx \\
    &= c_n \frac{1}{|Q|^\frac{1}{n'}} \int_Q |\nabla f(x)|dx.
\end{align*}
We now insert the weight in the previous estimate. Since
\begin{equation*}
    \frac{w(Q)}{|Q|} \le \inf_{x\in Q} M(w\chi_Q)(x),
\end{equation*}
we have
\begin{align*}
    \left(\tau_{\frac{t_1}{2}}(|f-f_Q|)\right)_Q
    &\le c_n \frac{1}{w(Q)^\frac{1}{n'}} \inf_{x\in Q} M(w\chi_Q)(x)^\frac{1}{n'} \int_Q |\nabla f(x)|dx \\
    &\le c_n \frac{1}{w(Q)^\frac{1}{n'}} \int_Q |\nabla f(x)|M(w\chi_Q)(x)^\frac{1}{n'}dx \\
    &= c_n w(Q)^\frac{1}{n} \left\| \left| \nabla f \right| \frac{M(w\chi_Q)^\frac{1}{n'}}{w} \right\|_{\strt{2.3ex} L^1\left(Q,\frac{w(x)dx}{w(Q)}\right)} \\
    &\le c_{n,p} w(Q)^\frac{1}{n} \left\| \left| \nabla f \right| \frac{M(w\chi_Q)^\frac{1}{n'}}{w} \right\|_{\strt{2.3ex} L^{p,\infty}\left(Q,\frac{w(x)dx}{w(Q)}\right)} \\
    &= c_{n,p}A.
\end{align*}
In the last inequality, we used Kolmogorov's inequality again. Choosing the constant $C_{n,p} = 4c_{n,p}$ in the definition of $t_0$, and recalling that $t_1\ge t_0$, we obtain
\begin{equation*}
    \left(\tau_{\frac{t_1}{2}}(|f-f_Q|)\right)_Q \le \frac{t_1}{4}.
\end{equation*}
As a result,
\begin{equation*}
    E_{t_1} \subset \left\{x\in Q:\left| \tau_{\frac{t_1}{2}}(|f-f_Q|)(x)-\left(\tau_{\frac{t_1}{2}}(|f-f_Q|)\right)_Q \right| \ge \frac{t_1}{4}\right\}.
\end{equation*}

Applying Theorem \ref{Thm FPW} to $\tau_{\frac{t_1}{2}}(|f-f_Q|)$ with the measure $w(x)dx$, we get
\begin{align*}
    \frac{t_1}{4}w(E_{t_1})^\frac{1}{n'}
    &\le \frac{t_1}{4} w\left( \left\{ x\in Q: \left| \tau_{\frac{t_1}{2}}(|f-f_Q|)(x) - \left(\tau_{\frac{t_1}{2}}(|f-f_Q|)\right)_Q \right| \ge \frac{t_1}{4} \right\} \right)^\frac{1}{n'} \\
    &\le \left\| \tau_{\frac{t_1}{2}}(|f-f_Q|) - \left(\tau_{\frac{t_1}{2}}(|f-f_Q|)\right)_Q \right\|_{\strt{2.3ex} L^{n',\infty}(Q, w)} \\
    &\le c_n \int_Q \left|\nabla \big( \tau_{\frac{t_1}{2}}(|f-f_Q|)\big)(x)\right| M^c(w\chi_Q)(x)^\frac{1}{n'}dx \\
    &\le c_n \int_{\left\{x\in Q:\frac{t_1}{2}<|f(x)-f_Q|<t_1\right\}} |\nabla f(x)| \frac{M(w\chi_Q)(x)^\frac{1}{n'}}{w(x)} w(x)dx.
\end{align*}
Using the Lorentz-space version of Hölder's inequality \eqref{Holder Lorentz}, it follows that, 
\begin{equation*}
    \frac{t_1}{4}w(E_{t_1})^\frac{1}{n'} \le c_{n,p} \left\| \chi_{\left\{x\in Q:\frac{t_1}{2}<|f(x)-f_Q|<t_1\right\}} \right\|_{\strt{2.3ex}L^{p',1}(Q,w)} \left\| \left| \nabla f \right| \frac{M(w\chi_Q)^\frac{1}{n'}}{w} \right\|_{\strt{2.3ex}L^{p,\infty}(Q,w)}.
\end{equation*}
Moreover, by \eqref{level comparison fQ}, 
\begin{align*}
    \left\| \chi_{\left\{x\in Q:\frac{t_1}{2}<|f(x)-f_Q|<t_1\right\}} \right\|_{\strt{2.3ex} L^{p',1}(Q,w)}
    &= c_p w\left(\left\{x\in Q:\frac{t_1}{2}<|f(x)-f_Q|<t_1\right\}\right)^\frac{1}{p'} \\
    &\le c_{n,p} w\left(E_{t_1}\right)^\frac{1}{p'}.
\end{align*}
Therefore,
\begin{equation*}
    t_1 w(E_{t_1})^\frac{1}{n'} \le c_{n,p} w(E_{t_1})^\frac{1}{p'} \left\| \left| \nabla f \right| \frac{M(w\chi_Q)^\frac{1}{n'}}{w} \right\|_{\strt{2.3ex} L^{p,\infty}(Q,w)}.
\end{equation*}
Since $\frac{1}{n'}-\frac{1}{p'}=\frac{1}{p^*}$, we obtain
\begin{equation}\label{key level estimate fQ}
    t_1 w(E_{t_1})^\frac{1}{p^*} \le c_{n,p} \left\| \left| \nabla f \right| \frac{M(w\chi_Q)^\frac{1}{n'}}{w} \right\|_{\strt{2.3ex} L^{p,\infty}(Q,w)}.
\end{equation}
Combining everything together, we have
\begin{align*}
    \left\| f-f_Q\right\|_{L^{p^*,\infty}\left(Q,\frac{w(x)dx}{w(Q)}\right)}
    &\le 2t_1 \left(\frac{w(E_{t_1})}{w(Q)}\right)^\frac{1}{p^*} \\
    &\le c_{n,p} \left(\frac{1}{w(Q)}\right)^\frac{1}{p^*} \left\| \left| \nabla f \right| \frac{M(w\chi_Q)^\frac{1}{n'}}{w} \right\|_{\strt{2.3ex} L^{p,\infty}(Q,w)} \\
    &= c_{n,p} w(Q)^\frac{1}{n} \left\| \left| \nabla f \right| \frac{M(w\chi_Q)^\frac{1}{n'}}{w} \right\|_{\strt{2.3ex} L^{p,\infty}\left(Q,\frac{w(x)dx}{w(Q)}\right)} \\
    &= c_{n,p}A.
\end{align*}
This proves \eqref{Thm PS 4 normalized Eq}, and concludes the proof. 
\end{proof}

\end{document}
